% Part: set-theory
% Chapter: story
% Section: urelements

\documentclass[../../../include/open-logic-section]{subfiles}

\begin{document}

\olfileid{sth}{story}{urelements}
\olsection{উর-উপাদান থাকবে কি?}

পরের কয়েকটি অধ্যায়ে আমরা সেটের সঞ্চয়ী-পর্যায়ক্রমিক ধারণা থেকে
স্বতঃসিদ্ধগুলি বের করার চেষ্টা করব। কিন্তু তার আগে
\emph{উর-উপাদান} সম্পর্কে আরও কিছু বলা দরকার।

\olref[approach]{sec}-এর ছবিতে কেবল পুরোনো \emph{সেট} থেকে
নতুন সেট গঠন করা যায়। তবে আমরা চাইতে পারি যেন কিছু
\emph{অ-সেট}—গরু, শুয়োর, বালুকণা কিংবা অন্য কিছু—সেটের
!!{element}s হতে পারে। সে ক্ষেত্রে শুরুতে কয়েকটি মৌলিক
উপাদান, অর্থাৎ \emph{উর-উপাদান}, থাকবে। তারপর বলব,
প্রত্যেক পর্যায়~$S$-এ উর-উপাদান অথবা $S$-এর আগের
পর্যায়গুলিতে গঠিত সেট—যেকোনো মিশ্রণে—দিয়ে তৈরি
``সম্ভব সব'' সংগ্রহকে সেট হিসেবে গঠন করা হয়। তখন ছবিটি
বরং এ রকম হবে:
\begin{center}
  \begin{tikzpicture}[scale=0.6]
  \tikzset{cut_here/.style={densely dotted}}
  \draw (-4,6) -- (-1, 0) -- (1,0) -- (4,6);
  \draw[cut_here] (-5,8)--(-4,6);
  \draw[cut_here] (5,8)--(4,6);
  \draw(-1.5, 1)--(1.5, 1);
  \draw(-2, 2)--(2, 2);
  \draw(-2.5, 3)--(2.5,3);
  \draw(-3, 4)--(3, 4);
  \draw(-3.5, 5)--(3.5, 5);
  \draw(-4, 6)--(4, 6);
  \node[label] at (2, 0) {\small 0};
  \node[label] at (2.5, 1) {\small 1};
  \node[label] at (3, 2) {\small 2};
  \node[label] at (3.5, 3) {\small 3};
  \node[label] at (4, 4) {\small 4};
  \node[label] at (4.5, 5) {\small 5};
  \node[label] at (5, 6) {\small 6};
  \end{tikzpicture}
\end{center}
এখন আমাদের সিদ্ধান্ত নিতে হবে: \emph{উর-উপাদান অনুমোদন করব কি?}

দার্শনিক দিক থেকে উর-উপাদানকে তত্ত্বে রাখা যুক্তিসঙ্গত।
মূল কারণ হলো সেটতত্ত্বকে \emph{প্রয়োগযোগ্য} করা। বিষয়টি
বোঝাতে \olref[sfr][siz][]{chap}-এর কথা মনে করি: দুটি সেট
$A$ ও~$B$ সমান আকারের, অর্থাৎ $\cardeq{A}{B}$, যদি এবং
কেবল যদি তাদের মধ্যে একটি একৈক-সমাপতিত অপেক্ষক থাকে।
এখন মাঠের গরুগুলি এবং খোঁয়াড়ের শুয়োরগুলি যদি দুটি সেট
গঠন করে, তাহলে ``গরু ও শুয়োর সমসংখ্যক'' দাবিটিকে
সেটতাত্ত্বিকভাবে ব্যাখ্যা করা যায়। কিন্তু উর-উপাদান নিষিদ্ধ
করলে গরু ও শুয়োরেরা \emph{সেট গঠন করে না}; তখন ওই
সেটতাত্ত্বিক ব্যাখ্যা আর পাওয়া যাবে না। আসলে সেট ছাড়া
অন্য কিছুর ক্ষেত্রে সেটতত্ত্ব প্রয়োগের কোনো সরল উপায়ই
আমাদের থাকবে না। (উর-উপাদান রাখার আরও কারণের জন্য
\citealt[pp.~vi, 24, 50--1]{Potter2004} দেখুন।)
% BN-SRC-779 TARGET-CORRECTION-BEGIN
\emph{পাঠ-স্পষ্টীকরণ।} এখানে বাইরের বস্তুগুলিকে সরাসরি সদস্য
করার কথা বলা হচ্ছে। বিশুদ্ধ সেটতত্ত্বেও তাদের সেট দিয়ে সংকেতায়িত
করে এবং সেই প্রতিনিধিত্ব ব্যাখ্যা করে গণনা ও একৈক-সমাপতন
প্রয়োগ করা যায়। উর-উপাদান অনুমোদন করলে সদস্যভিত্তিক সমতার
স্বতঃসিদ্ধে দুই তুলনীয় বস্তুকে সেট বলে সীমিত করতে হবে: ভিন্ন
উর-উপাদান এবং শূন্য সেট—সকলেরই কোনো সদস্য নেই বলে তাদের
এক করে ফেলা চলে না। বইটির পরের বিশুদ্ধ সেটতত্ত্বে উর-উপাদান
অনুমোদিত নয়; সেই প্রকার-সীমার প্রয়োজন সেখানে পড়ে না।
% BN-SRC-779 TARGET-CORRECTION-END

গাণিতিক চর্চায় অবশ্য উর-উপাদান অনুমোদন করা বেশ বিরল।
এর আংশিক কারণ, উর-উপাদান ছাড়া সেটতত্ত্ব সাজানো
\emph{সামান্য} সহজ। তবে কখনও উর-উপাদান বাদ দেওয়ার
আরও তাৎপর্যপূর্ণ যুক্তিও পাওয়া যায়:
\begin{quote}
  সেটতত্ত্ব গণিতের ভিত্তি—এই বিশ্বাস অনুযায়ী শুধু সেটের
  কথা বলেই আমাদের সমস্ত গণিত প্রকাশ করতে পারা উচিত। তাই
  আমাদের চলরাশির বিস্তারে গরু-শুয়োরের মতো বস্তু থাকা উচিত
  নয়।
  %But if $C$ is a cow, $\{C\}$ is a set, but not a legitimate mathematical object.
  \citep[p.~8]{Kunen1980}
\end{quote}
অতএব প্রয়োগযোগ্যতায় জোর দিলে উর-উপাদান \emph{রাখার}
পক্ষে যুক্তি মেলে; ভিত্তি হিসেবে গণিতকে বিশুদ্ধ সেটতত্ত্বে
নামিয়ে আনার লক্ষ্য হলে সেগুলি \emph{বাদ দেওয়ার} পক্ষে
যুক্তি মিলতে পারে। একটু অলসতাও বাদ দেওয়ার পক্ষেই যায়।

আমরা সবচেয়ে সহজ পথটি নেব। তবে এর একটি শিক্ষণগত
যুক্তিও আছে। আমাদের লক্ষ্য সেটতত্ত্বের সেই মূল বিষয়গুলি
পাঠককে শেখানো, যেগুলি জানলে সেটতত্ত্বের দর্শনে প্রবেশ
করা যায়। সেই দর্শনের অধিকাংশ লেখাতেই \emph{উর-উপাদানবিহীন}
সেটতত্ত্ব আলোচনা করা হয়। তাই বইটিও ওই সাহিত্যের
কাছাকাছি থাকলে হয়তো বেশি কাজে লাগবে।

\end{document}
